\documentclass[11pt,a4paper]{article}
\usepackage[margin=25mm]{geometry}
\usepackage{amsmath,amssymb,amsthm,hyperref}
\newtheorem{lemma}{Lemma}
\newcommand{\gen}[1]{\langle#1\rangle}
\title{Localisation in the paired-cut equality case}
\author{Human Compression project: L044 working continuation}
\date{}
\begin{document}\maketitle
Use the notation and square-normal-form theorem of \emph{Paired orbit
cuts in string actions}. Suppose $2r=n$. Among the paired selections
whose only cycles are alternating squares let $Q$ minimise the number
$q$ of squares.

\begin{lemma}\label{all}
Every such minimum-square selection has exactly one square in every
component.
\end{lemma}
\begin{proof}
If some minimum-square selection had two squares in one component,
choose, over all such selections, a pair at the smallest distance.
The incident-square, ordinary-transport and global middle-colour
moves in the square-normal-form proof contradict first the minimal
square count and then this minimal distance. Thus every component
has at most one square. The exact defect $2r=n-c+q$ and $2r=n$ give
$q=c$, so there are no tree components.
\end{proof}

\begin{lemma}[Localising cycle colours]\label{colour}
Let $S$ be the square in a component $C$, and let $a$ be one of its
two colours. Every eligible $a$-edge of the full permutation graph
has both endpoints in $C$.
\end{lemma}
\begin{proof}
If both endpoints lay in another component, its connecting path in
$Q$ would avoid $a$, contradicting eligibility. If the endpoints
lay in different components, delete one selected $a$-edge of $S$
and insert the new edge. The deletion destroys $S$ without
disconnecting $C$, while the new edge joins different components.
No cycle is created, so the number of squares decreases. This also
is impossible. The only remaining location is within $C$.
\end{proof}

\begin{lemma}[Full neighbours of square vertices]\label{neighbour}
Every edge of the full permutation graph incident with a vertex of
$S$ has its other endpoint in $C$.
\end{lemma}
\begin{proof}
Let the square colours be $a,b$, and let $v\in S$. The cases of
colours $a,b$ are internal to $S$. If a colour $k$ commutes with
one of them, say $a$, then the image of an $a$-edge under
$\rho_k\in G_a$ is an actual eligible $a$-edge. Lemma~\ref{colour}
places both its endpoints, including $v^{\rho_k}$, in $C$.

The remaining case has $a=i-1,b=i+1,k=i$. Use the global
middle-colour exchange, obtaining $Q'$ with the same minimum $q$.
Viewed as a graph operation on $Q^{\rho_i}$, this removes
$S^{\rho_i}$ and inserts $S$. Removing the four edges of
$S^{\rho_i}$ splits $C^{\rho_i}$ into four trees; every other
component remains unicyclic. All four vertices of the inserted
$S$ must lie within $C^{\rho_i}$. Otherwise the component containing
that vertex in $Q'$ would contain both the new square $S$ and the
unchanged old square of another component, contrary to
Lemma~\ref{all}. Therefore $S\subseteq C^{\rho_i}$, equivalently
$S^{\rho_i}\subseteq C$. In particular $v^{\rho_i}\in C$.
\end{proof}

\begin{lemma}[Ordinary slides preserve the component partition]
Let a selected bridge $v\mathrel{-_\ell}w$ leave $S$, with $\ell$
commuting with a square colour $a$. The ordinary square slide may
be made without changing any component's vertex set.
\end{lemma}
\begin{proof}
Write $v'=v^{\rho_a}$ and $w'=w^{\rho_a}=(v')^{\rho_\ell}$.
By Lemma~\ref{colour}, the eligible new edge $ww'$ lies within $C$.
Delete the opposite $a$-edge of $S$, add $ww'$, and complete the
parallel $\ell$-edge $v'w'$ if necessary, deleting the other
selected $\ell$-edge. As in ordinary transport, the number of
squares cannot increase and no other cycles are created; minimality
forces the square count to stay $q$.

If the removed $\ell$-edge were outside $C$, it would be a bridge:
every cycle containing it would also require the selected
$\ell$-edge $vw$, which lies in $C$. Removing that outside bridge
leaves a tree component, and all inserted edges are within $C$.
This contradicts Lemma~\ref{all}. Thus every changed edge lies in
$C$. Its new selection has one square and cannot leave a tree
component, again by Lemma~\ref{all}; hence it remains connected.
Every other component is untouched.
\end{proof}

\begin{lemma}[The transported partition]\label{partition}
For a square $S$ of colours $i-1,i+1$, the global middle-colour
exchange has component partition exactly the image under $\rho_i$
of the old component partition.
\end{lemma}
\begin{proof}
The proof of Lemma~\ref{neighbour} shows that every vertex of the
inserted $S$ lies in $C^{\rho_i}$. Deleting $S^{\rho_i}$ split this
component into four trees. Inserting $S$ must reconnect all four:
any omitted piece would remain a tree, contrary to Lemma~\ref{all}.
All other components are the untouched images of the old components.
Thus the new component containing $S$ has exactly the vertex set
$C^{\rho_i}$, as claimed. The claim is transport of the partition,
not invariance of the old partition.
\end{proof}

\begin{lemma}[Connected equality normal form]\label{connected}
If $2r=n$, a minimum-square paired selection is connected and has
exactly one square.
\end{lemma}
\begin{proof}
The minimum number $q$ of squares is fixed. If a minimum-square
selection is connected then Lemma~\ref{all} already gives $q=1$;
the same is then true of every minimum-square selection. Suppose
instead that all such selections are disconnected. The full
permutation graph is connected by transitivity. Hence for any of
these selections there is a full graph edge between two components,
say $C,D$, and there are selected paths from their unique squares
$S,T$ to the respective endpoints. Choose, over \emph{all}
minimum-square selections and \emph{all} such edges and selected
tails, a connector $P$ with the fewest edges. Its only unselected
edge is this cross-edge, whose colour will be denoted by $x$.

Both selected tails are nonempty, by Lemma~\ref{neighbour}. Suppose
the first edge from $S$ has colour $\ell$ and commutes with a colour
$a$ of $S$. Apply the ordinary slide, with notation $v,v',w,w'$
as above. If its other selected $\ell$-edge lay on $P$, then $w'$
would lie outside $C$: on the $w$ side of the bridge $vw$ a path
to $w'$ would avoid $a$, whereas on its other side a path from
$v'$ to $w'$ would avoid $\ell$. The same argument works whether
the other $\ell$-edge lies later in the $C$ tail or in the $D$
tail. Both alternatives contradict eligibility, exactly as in the
repeated-colour case of ordinary transport. But $ww'$ is an
eligible $a$-edge, so Lemma~\ref{colour} requires $w'\in C$.
Thus the other $\ell$-edge is not on $P$.

Consequently the ordinary slide preserves the remaining connector.
It also preserves its two component vertex sets by the preceding
slide lemma. Its new square includes $w$, shortening the selected
tail by one, contrary to the choice of $P$. Thus the first edge
at $S$ must have a middle colour $i$ with square colours $i-1,i+1$.
The same argument at $T$ gives a middle colour $j$ and square
colours $j-1,j+1$.

The two square-colour sets are disjoint. In particular the
cross-edge colour $x$ lies outside at least one of them; choose
that endpoint, say $S$. Then $\rho_i$ commutes with $\rho_x$,
including when $x=i$. Perform the global middle-colour exchange
at $S$. Every selected edge of the connector other than the first
is outside $S$, so is transported to a selected edge. The single
full cross-edge is transported to an actual $x$-edge because of
this commutation. By Lemma~\ref{partition} it still crosses exactly
the two transported components. The transported tail begins at
$w^{\rho_i}=v\in S$ and ends at the transported $T$, and has one
less edge. It is therefore an admissible shorter connector, the
final contradiction.
\end{proof}

\paragraph{Attribution and remaining step.}
The source's Proposition~4.12 gives the relevant disconnected
reduction using three local connector configurations. The proof
above instead combines the root-supplied global middle-colour move
with minimum-square localisation. It makes no new source, priority
or conjecture claim. The connected equality selection still needs
the source's rigidity argument, giving a preserved system of pairs
for $n\geq9$ and its exact two ladder groups. Connectedness alone
does not yet exclude the alternating action.
\end{document}
